\documentclass[10pt]{article}
\usepackage{template/template}
\usepackage{tikz}
\NewsSetKeys{date={2026年11月13日},vol={031},title={每日有趣新闻},author={编辑部},note={待补编短刊。选材以11月11—12日为主；数学栏目另收11月9日新预印本，明确标注日期。尚未通过正式发行验收。}}
\begin{document}
\begin{documents}
\newsitem{能源供应链的预订泡沫：AI数据中心究竟需要多少燃气轮机？}{科技与经济｜11月11日｜深度观察}{
AI基础设施的建设速度不仅取决于芯片，还取决于能够实际交付的电力。路透社Breakingviews在11月11日的评论中讨论软银取消美国数据中心设备订单并确认约37亿美元费用的事件。报道指出，软银此前预订了超过13吉瓦燃气轮机产能，部分规模超出现有主要项目需求。供给紧张时，提前预订能够保留未来建设选项；但如果多个项目都以预留设备作为可行性证明，名义订单就可能超过最终投入运行的容量。\cite{softbank}
这里必须区别三种数字：供应商订单簿代表合同或预订，项目公布的电力规模代表规划，真正可持续供电的算力中心则必须完成设备交付、施工、并网、许可和运营。三者不能直接相加或互换。订单取消说明个别买方调整资本配置，并不等于AI计算需求突然消失；反过来，设备订单持续增长也不能直接证明每座计划中的数据中心都能盈利。此文是金融评论而非独立审计，具体订单和财务数字仍应以企业披露进一步核对。
从技术角度看，燃气轮机可以提供相对稳定的现场电源，却涉及燃料、维护、排放与长期锁定成本。即使机组本身交付，供电系统还需要变压、冷却、冗余与调度；用电能力也不等于GPU利用率。市场最值得观察的不是某一次取消订单，而是供应商实际交货、客户资本开支、数据中心并网容量和推理服务收入能否在多个季度保持一致。
}
\newsitem{新一轮对俄能源制裁：制裁清单不等于出口量立即归零}{国际政治与能源｜11月12日}{
英国政府宣布将俄罗斯石油生产商Slavneft及一批能源相关实体、船舶纳入制裁。路透社报道，英国方面称此举使其覆盖俄罗斯十大石油生产商；其政治目标是限制战争收入，并向俄罗斯施加谈判压力。\cite{ukrussia} 这属于英国政府的政策表述，而非对制裁效果的实证结论。制裁能否改变贸易流向，取决于司法管辖、保险、运输、第三国购买与执行能力。短期价格变化还会受到中东冲突、库存与市场预期影响。应避免将新增名单直接等同于俄方石油出口立刻减少同等规模。
}
\newsitem{数学拾趣：不用复Monge--Ampère方程，能否重新证明Calabi--Yau定理？}{数学｜11月9日预印本，非当日新闻}{
Bennett Chow与Yuan Liao在arXiv:2611.12800提出紧致Kähler流形上Calabi--Yau定理的一个几何证明：在实第一陈类为零时，每个Kähler类包含唯一Ricci平坦度量，并证明Kähler--Ricci流光滑收敛至该度量。\cite{chowliao} 简单说，Kähler类固定了某种整体几何约束，而Ricci平坦要求曲率的一个重要收缩量为零。困难不在于写出流方程，而在于防止长期演化中几何尺度退化，并取得足够强的估计以证明极限存在且光滑。
作者在摘要中说明，利用平均点对距离的控制、相对体积比较与直径估计，再以热方程梯度估计及其他内部估计推出指数收敛。其特色是声称不调用原有Calabi--Yau存在性定理，也不依赖标量复Monge--Ampère方程的全局估计。这是新证明路线的主张，而不是推翻原定理。本文为11月9日提交的v1预印本；本刊核对了原始摘要，尚未逐步审查证明，也未核实同行评议接受状态，不能宣称证明已经获得学界确认。
}
\newsitem{AI生成研究成果，机器验过的究竟是什么？}{AI与数学方法论｜11月9日预印本}{
Eigner、Süli和Hollós提出一个让人类研究者参与审计的工作流，强调自动证明器的通过结果只对被检查的形式对象负责，不能自动保证形式陈述忠实表达原始自然语言命题。作者展示了一个含糊主张分裂为两种解释、其中一种被反驳而另一种被证明的案例。\cite{audit} 这是作者报告的实验，不能据此推断所有自动形式化系统均有相同比例的语义错误。论文仍为预印本；文章的主要启示是把命题忠实性、工具可信性和最终责任分开。
}
\newsitem{国际局势速记：谈判进展与战争结束不是同义词}{外交与战争｜11月12日}{
《卫报》11月12日的滚动报道记述，美俄代表会晤后宣布拟继续与乌克兰及欧洲方面磋商；与此同时，俄方对领土和外国维和力量的立场仍与乌克兰及部分欧洲国家存在显著冲突。\cite{ukrainetalk} 滚动报道汇集多方声明，应分别识别各方立场。安排下一轮谈判是程序性进展，不构成停火或和平协议已经签署的证据。
}
\end{documents}
\bibliographystyle{plain}
\bibliography{refs}
\end{document}